\SourceLeaf{543}
\pdfbookmark[0]{Magnetnadel in Kopenhagen und Mailand. 1834}{copenhagen_mailand_1834}
\begin{center}\itshape
Beobachtungen der Schwankungen der Magnetnadel in Kopenhagen und Mailand\\
am 5. und 6. November 1834.
\end{center}
\begin{center}
\hrule\vspace{2mm}
\textsc{Schumacher}. Astronomische Nachrichten Nr.\ 276. 21. März 1835.
\vspace{2mm}\hrule\vspace{5mm}
\end{center}
Seit der Fertigstellung des hiesigen magnetischen Observatoriums werden hier unter anderem an bestimmten, im Voraus festgelegten Tagen die Schwankungen der magnetischen Deklination regelmäßig in kurzen Zeitabständen ohne Unterbrechung beobachtet. Anfangs wurden dafür dieselben Termine gewählt, die Herr \textsc{von Humboldt} schon vor mehreren Jahren angesetzt hatte. Seit dem vergangenen Frühjahr haben sich bereits recht viele Astronomen und Physiker Apparate beschafft, die dem hiesigen ähneln, den ich an anderer Stelle hinreichend beschrieben habe. Sie nehmen an den vereinbarten Beobachtungen teil. Schon die ersten auf diese Weise gewonnenen gleichzeitigen Beobachtungen am 4. und 5. Mai in Göttingen und Waltershausen (einem Gut in der Gegend von Schweinfurt, wo Herr \textsc{Sartorius} mit einem kleinen, sonst aber dem hiesigen ganz ähnlichen Apparat beobachtete) zeigten eine außerordentlich bemerkenswerte Übereinstimmung im vielfach hin und her springenden Verlauf der Schwankungen, bei den kleineren ebenso wie bei den größeren. Ähnliche Ergebnisse haben sich seitdem bei späteren Terminen mit Beteiligung von Leipzig, Berlin, Braunschweig und Kopenhagen vielfach wiederholt. Einige Beispiele sind in \textsc{Poggendorffs} Annalen als grafische Darstellungen veröffentlicht worden.

Je weiter sich die Teilnahme an diesen vereinbarten Beobachtungen ausbreitet, desto interessantere und fruchtbarere Ergebnisse sind natürlich zu erwarten. Ich wiederhole daher hier die bereits an anderer Stelle veröffentlichte Mitteilung. Seit sich so klar gezeigt hat, dass in \textit{sehr} kurzen Zeitabständen beobachtet werden muss, haben wir uns veranlasst gesehen, die Termine zu ändern. Wir haben ihre Anzahl von 8 auf 6 im Jahr und ihre Dauer von 44 auf 24 Stunden verringert. Derzeit ist der letzte Samstag jedes ungeraden Monats festgelegt, von Göttinger Mittag bis zum Mittag des folgenden Tages. Zu jedem dieser Haupttermine kommen zwei Nebentermine hinzu, nämlich am
\SourceLeaf{544}darauffolgenden Dienstag und Mittwoch abends von 8 bis 10 Uhr. Ausführlichere Angaben, auch zur Beobachtungsweise, finden sich in \textsc{Poggendorffs} Annalen, Band~33. [S.\ 528 dieses Bandes.]

Das bemerkenswerteste bisher erzielte Ergebnis liefern die gleichzeitigen Beobachtungen in Kopenhagen und Mailand am 5. und 6. November dieses Jahres. Der Termin gehörte noch zur früheren Regelung, von deren Änderung die Beobachter an beiden Orten noch keine Nachricht erhalten hatten. In Kopenhagen, wo jetzt unter der Leitung von Herrn Staatsrat \textsc{Oersted} ein dem hiesigen ganz ähnliches magnetisches Observatorium eingerichtet ist, wurde eine Nadel von derselben Stärke wie die hiesige verwendet (vier Pfund schwer). In Mailand beobachteten Herr \textsc{Sartorius} und Dr.\ \textsc{Listing} an der dortigen Sternwarte mit der oben erwähnten kleineren Nadel. Sie wurden dabei von Herrn \textsc{Kreil}, einem Schüler der Sternwarte, unterstützt. Ich gestehe, dass mich trotz der vielen früheren Erfahrungen ähnlicher Art das Ausmaß der Übereinstimmung an zwei mehr als 150 Meilen voneinander entfernten Orten überrascht hat. Schon der bloße Anblick der beigefügten grafischen Darstellung spricht für sich. Ich füge ihr nur einige Erläuterungen und Bemerkungen bei.

Da mir der Wert der Skalenteile in Kopenhagen anfangs noch unbekannt war, entwarf ich die Zeichnung in einem Maßstab, in dem die Anomalien ungefähr gleich groß erscheinen. Dazu ließ ich die Seite jedes Gitterquadrats neun Skalenteilen der Kopenhagener und drei Skalenteilen der Mailänder Beobachtungen entsprechen. Ein Skalenteil in Kopenhagen beträgt übrigens $21''576$, einer in Mailand $29''341$. In Winkelmaßen waren die Kopenhagener Bewegungen also etwa $2{,}2$-mal so groß wie die Mailänder. Will man daraus auf das Verhältnis der wirkenden Kräfte schließen, muss man berücksichtigen, dass diese Erscheinungen nur als Störungen der horizontalen erdmagnetischen Kraft an beiden Orten zu betrachten sind. Dieselbe störende Kraft muss an einem Ort, an dem die horizontale Kraft kleiner ist, größere Änderungen hervorrufen als an einem anderen, an dem sie größer ist. Nach \textsc{Hansteens} schöner Karte im 7. Band der Astronomischen Nachrichten\HistoricalNote{544} schätzte ich das Verhältnis der horizontalen Kraft des Erdmagnetismus in Kopenhagen und Mailand auf
\SourceLeaf{545}1 zu $1{,}23$. Danach ließe sich das Verhältnis der störenden Kräfte, die an jenen Tagen die beträchtlichsten Anomalien in Kopenhagen und Mailand hervorgerufen haben, auf etwa $1{,}8$ zu 1 schätzen. Wie viel besser werden wir künftig über solche rätselhaften Naturkräfte urteilen können, wenn uns erst ähnliche gleichzeitige Beobachtungen von vielen weit voneinander entfernten Orten vorliegen.

Neben der überraschend großen Übereinstimmung im Verlauf der Anomalien bemerken wir allerdings auch Unterschiede. Erstaunlich scheint dabei, wie klein diese Unterschiede im Vergleich sind. Wir kennen die Ursachen der Erscheinungen freilich noch gar nicht. Gerade bei ihrem vielfältigen Wechsel erscheint es unnatürlich, anzunehmen, dass sie alle von einem einzigen Punkt aus wirkten. Einige Ursachen mögen hier, andere dort ihren Sitz gehabt haben. So mögen während der 44 Stunden auch manche Kräfte aus ganz anderen Gegenden mitgewirkt haben, die an den beiden Orten in einem ganz anderen Verhältnis standen. Dass die Kurve für Mailand im Allgemeinen viel unruhiger aussieht als die für Kopenhagen, erklärt sich übrigens von selbst daraus, dass in Mailand alle 5 Minuten, in Kopenhagen alle 10 Minuten beobachtet wurde. Bei den längeren Zwischenzeiten mussten folglich manche kleineren und schneller wechselnden Anomalien unbemerkt bleiben.

Auch wenn das Interesse an diesen Forschungen keiner Verstärkung bedarf, glaube ich doch, noch einen Umstand hervorheben zu müssen, der die Astronomen besonders angeht. Ob die bei diesen Bewegungen wirkenden Kräfte eine messbare Zeit brauchen, um sich über große Entfernungen fortzupflanzen, wissen wir noch nicht. Diese interessante Frage wird aber ohne Zweifel in
\SourceLeaf{546}Zukunft beantwortet werden. Ist die Zeit unmessbar klein, können solche Beobachtungen rasch ansteigender und abfallender Bewegungen zur Bestimmung geografischer Längen dienen. Unter günstigen Umständen können sich diese Längenbestimmungen sogar mit den genaueren messen. Aus beobachteten Bewegungen in Göttingen und Leipzig habe ich unter jener Voraussetzung schon mehrmals den Längenunterschied bis auf eine halbe Zeitminute richtig ableiten können. Zuweilen treten jedoch so schnelle Bewegungen auf, dass sich daraus noch viel genauere Zeitbestimmungen ableiten lassen. Die stärksten Bewegungen, die mir bisher begegnet sind, fanden am 7. Februar dieses Jahres statt, als die Nadel den ganzen Tag überaus unruhig war. Ich beobachtete Bewegungen von 17 Skalenteilen oder 6 Bogenminuten in einer einzigen Zeitminute. Sie hielten einige Minuten regelmäßig an, wurden dann nach und nach langsamer und gingen anschließend in die entgegengesetzte Richtung über. Solche Erscheinungen könnten, wenn man sie an zwei Orten mit guten Apparaten sorgfältig verfolgt (die selbst bei einzelnen Beobachtungen eine Genauigkeit von wenigen Bogensekunden liefern), den Längenunterschied bis auf eine Zeitsekunde genau ergeben, sofern die Kräfte in unmessbar kleiner Zeit wirken. Jedenfalls ist klar, wie wichtig es zur Klärung der Frage sein wird, dass alle Beobachter, denen die Mittel dazu zur Verfügung stehen, stets auf eine gute Zeitbestimmung achten.

Schließlich bemerke ich noch, dass die Beobachter in Mailand die dortige Inklination am 2. November mit einem \textsc{Lenoir}schen Inklinometer $=63^0\,55'\,26''$ gefunden haben.

